\chapter{可信执行、身份与证明}
\label{chap:trusted-execution-identity}

可信启动只能证明系统加载了经过授权的软件，不能自动保护运行时密钥、证明设备当前状态或隔离高价值操作。完整信任体系还需要硬件唯一根、受保护执行环境、密钥派生、设备证明和生命周期管理。

这些机制的工程价值在于建立可验证边界：普通软件不能导出哪些密钥，哪些状态不能回放，谁能根据什么证据信任一台设备，以及信任失效后如何恢复。产品安全需求、漏洞响应、生产与退役流程见 \mbox{第~\ref{chap:security-lifecycle}~章}；本章聚焦信任机制、数据契约和验证方法。

\section{信任根与身份层次}

设备身份不应等同于可修改的序列号或 MAC 地址。根身份通常来自不可变密钥、PUF、eFuse/OTP、安全元件或在受控生产阶段注入的设备密钥。

一台产品通常同时具有多层身份：芯片根身份表示不可替换硬件，制造身份绑定板卡和产品型号，运营身份用于云端或企业网络，租户与用户身份表示所有权，会话身份则是短期且可撤销的访问凭据。层次之间应有明确背书和撤权关系，不能用一个永久序列号承担全部授权。

直接长期使用硬件根密钥会扩大泄漏影响。更合理的层次是：根密钥只参与派生，按用途生成启动验证、存储加密、设备认证和会话密钥。派生上下文应绑定设备、用途、算法版本和生命周期状态。

派生上下文本身就是长期格式。修改 purpose 编码、算法、设备标识或 measurement 含义，都可能导致旧数据永久无法解密。因此上下文应带 schema version，并为密钥轮换、主板更换和安全版本提升保留迁移路径。

\noindent\begin{minipage}{\textwidth}
\begin{lstlisting}[language=C,caption={用途隔离的密钥派生上下文（示意）}]
struct key_context {
    uint32_t schema_version;
    uint32_t purpose;
    uint8_t  device_id[16];
    uint8_t  firmware_measurement[32];
    uint8_t  salt[32];
};

derived_key = hkdf(root_secret, context, sizeof(context));
\end{lstlisting}
\end{minipage}

同一个派生密钥不能同时用于 TLS、磁盘加密和更新签名验证。密钥用途分离可以限制单一协议缺陷的影响范围。

\section{可信计算基与组件边界}

可信计算基（TCB）包含其失效可以破坏安全目标的所有硬件、固件和软件。把业务代码放入 Secure World 并不会消除缺陷，只会把它移入更难调试、审计和更新的 TCB。工程目标应是让 TCB 尽量小，接口稳定，并使普通世界破坏后仍不能越过规定边界。

\begin{table}[H]
  \centering
  \caption{常见信任组件的工程边界}
  \begin{tabular}{p{30mm}p{38mm}p{38mm}p{30mm}}
    \hline
    组件 & 主要能力 & 不自动保证 & 常见风险 \\
    \hline
    Boot ROM/OTP & 最初验证与根密钥 & 上层协议与业务正确 & 不可更新缺陷 \\
    TEE & 隔离执行与密钥服务 & TA 输入与 DMA 隔离 & 共享内存与 TCB 膨胀 \\
    TPM/安全元件 & 受保护密钥与证明 & 主机已在正确使用密钥 & 授权策略与性能 \\
    RPMB & 认证读写与防重放计数 & 机密性、大容量与高频写 & 写计数与一致性 \\
    \hline
  \end{tabular}
\end{table}

边界评审要同时检查 CPU 访问、DMA、中断、调试器、缓存、共享内存和复位后状态。仅有 CPU 世界位或 PMP 配置，不表示所有总线主设备和早期启动阶段都受同一策略约束。

\section{Arm TrustZone 与 OP-TEE}

Arm TrustZone 把处理器执行划分为安全与非安全状态，安全监控器负责世界切换；OP-TEE 等可信执行环境在 Secure World 中运行受信任应用（TA），普通世界客户端通过驱动和共享内存请求服务\cite{optee-doc}。GlobalPlatform TEE 体系结构和 Client API 定义了常见角色、会话与参数语义，但具体平台的中断控制器、内存控制器和安全外设仍由 SoC 集成决定\cite{globalplatform-tee-system-architecture,globalplatform-tee-client-api}。

TEE 适合保存小型密钥、执行签名、验证受保护计数器和隔离设备身份操作，不适合把网络栈、媒体框架或整个业务搬入安全世界。普通世界失陷后，攻击者仍可高频调用合法命令、替换共享内存、取消请求或制造资源耗尽，因此安全服务必须实现自己的身份、授权、配额和状态机。

客户端、内核驱动、TEE Core 与 TA 之间的命令 ABI 应显式规定：

\begin{itemize}
  \item 协议版本、命令编号、参数方向和最大长度；
  \item session 建立、取消、超时、并发和幂等语义；
  \item 共享内存在校验后、复制前和使用期间是否可能被再次修改；
  \item TA 崩溃、TEE 重启、普通世界重启和整机复位时的错误映射；
  \item 密钥不存在、权限不足与内部故障是否需要统一外部错误粒度。
\end{itemize}

共享内存参数始终来自非可信世界。TA 不能先按长度分配资源、随后才验证上限，也不能在两次读取之间假定内容不变；必要时应把固定上限的数据复制到安全内存后再解析。对 scatter-gather、嵌套偏移和结构体版本还要检查整数溢出、重叠区间和 ABI 对齐。

世界位只约束规定的 CPU 和总线路径。DMA、调试模块、安全中断、缓存维护、推测执行、共享 LLC 和电源管理都可能跨越边界。平台必须证明非安全总线主设备不能访问安全 DRAM，复位后防火墙在加载非可信软件前已经生效，并对高价值密码操作评估时序、缓存和功耗侧信道。

\section{RISC-V 隔离基础}

RISC-V 可以使用 PMP 限制较低特权级对物理地址的访问，并通过 M-mode 固件建立 S-mode 与 U-mode 边界\cite{riscvpri}。PMP entry 数量、粒度、匹配方式、锁定行为和增强扩展由具体实现决定，不能假定所有 RISC-V SoC 具有相同安全能力。

可信执行方案还需要安全启动、可信固件、隔离中断、DMA 防火墙和安全存储。只有 PMP 而没有外设访问控制时，恶意 DMA 主设备仍可能破坏受保护内存。

M-mode 固件通常属于 TCB。它负责建立 PMP、委托异常与中断、管理定时器和执行环境调用，因此一个过宽的 SBI 扩展或错误的 trap 处理就可能绕过上层隔离。固件应在启动早期设置默认拒绝区域，只向较低特权级开放明确内存和 MMIO 窗口，并验证 warm reset、secondary hart 启动和 suspend/resume 不会恢复到宽松配置。

实现评审应形成每个 SoC 独立的能力清单，包括 PMP entry、锁定位、debug authentication、IOMMU 或总线防火墙、安全中断路由、不可变启动级、密钥存储与防回滚计数。若某型号缺少其中一项，产品策略应明确补偿控制或降低安全声明，不能用同一软件镜像名称掩盖硬件边界差异。

\section{TPM、安全元件与 RPMB}

TPM 提供受保护对象、PCR、策略授权、密封和证明，其命令与数据结构由 TPM 2.0 Library Specification 定义\cite{tcg-tpm2-library}。安全元件适合保存设备私钥并执行有限密码操作；eMMC RPMB 提供认证读写与单调写计数。它们解决的层次不同，不应混为一种“安全芯片”。

PCR 保存按顺序扩展的度量值，sealed object 可以把解封条件绑定到 PCR、命令、授权值或策略签名。PCR 相等只表示所选度量序列相等，不表示未被度量的外设、配置和运行时进程可信。TPM 也不知道主机为何请求一次签名，授权策略必须把密钥用途、调用者和外部业务状态关联起来。

安全元件和 TPM 的吞吐、对象槽位、session 数量、算法集合与掉电行为通常远小于主 CPU 密码库。调用层应限制并发、缓存非秘密公钥材料，并把 BUSY、RETRY、设备复位与永久故障区分开。为了性能把私钥静默导出到普通内存，会直接破坏原有信任边界。

高频、大容量数据不应直接写入 RPMB。常见做法是在 RPMB 保存小型根元数据、版本计数或数据分区密钥，再由普通存储保存经过认证加密的大数据。

\section{可信存储与新鲜性}

可信存储至少要分别声明机密性、完整性、新鲜性和掉电一致性。加密防止读取，MAC 或 AEAD 检测篡改，单调 generation 或受保护计数器检测旧副本回放，事务协议则保证任意复位点都能恢复到一个完整状态。这四项缺少任何一项，都不能笼统称为“安全存储”。

一种常见结构是普通分区保存经过认证加密的数据对象，RPMB 或安全元件只保存小型根记录。根记录可以包含 schema version、当前 generation、对象树根摘要、密钥版本和提交状态：

\noindent\begin{minipage}{\textwidth}
\begin{lstlisting}[language=C,caption={可信存储根记录（示意）}]
struct trusted_root_record {
    uint32_t schema_version;
    uint32_t key_version;
    uint64_t generation;
    uint8_t  root_digest[32];
    uint8_t  commit_state;
};
\end{lstlisting}
\end{minipage}

提交时应先写入可验证的新数据，再原子推进受保护根；恢复时只接受由当前根可达且认证成功的对象。若先推进计数器再写大数据，掉电可能使最新 generation 永远缺失；若先覆盖旧数据又没有版本化副本，则根记录未提交时也无法回退。底层介质、flush 和真实掉电语义见 \mbox{第~\ref{chap:storage-reliability}~章}。

单调计数器不是无限资源。设计应计算量产测试、正常业务、证书轮换、异常重试和预期寿命内的最坏写入次数，并为计数器耗尽、RPMB key 未注入或控制器永久失败定义有界降级。安全状态不能是无限重启，也不能在错误后退回不校验 generation 的普通文件。

\section{Measured Boot 与事件日志}

Verified Boot 决定“是否允许启动”，Measured Boot 记录“实际启动了什么”。每一级对下一级镜像和关键配置计算 measurement，并扩展到受保护寄存器或证明日志。

measurement 必须覆盖实际改变行为的内容，例如 Bootloader 配置、DTB、内核命令行、initramfs 和安全策略。只度量内核镜像无法表示完整启动状态。

PCR 一类寄存器通常按如下形式扩展，而不是直接覆盖：

\[
  PCR_{\mathrm{new}} = H(PCR_{\mathrm{old}} \parallel measurement)
\]

这个结构把顺序纳入最终值，但只看 PCR 无法解释是哪一个组件产生差异。事件日志应记录 measurement、组件标识、加载阶段、算法和可验证元数据；验证方重新计算扩展链并与受保护 PCR 对比。事件日志本身可以放在普通内存中，因为篡改会导致重算不匹配，但解析器仍要防止畸形长度和资源耗尽。

度量边界必须覆盖所有可选择启动路径，包括正常槽、恢复系统、工厂镜像、可选 DT overlay、协处理器固件和安全策略。若某组件在度量之后还能被普通软件修改，则 measurement 只描述旧副本。动态加载模块、容器和长期运行状态需要单独的运行时度量或工作负载身份机制，不能无限延伸“启动可信”的含义。

\section{证明证据数据模型}

远程证明把身份、状态 claims、随机 challenge 和签名组合，使验证方判断设备是否处于可接受状态。证明不能只返回固件版本字符串，因为普通软件可以伪造该字符串。Arm PSA Attestation 等 profile 给出了面向受限设备的 claim 与 token 约定\cite{arm-psa-attestation-api}。

Evidence 应有明确 profile 和 schema version，并使用确定性编码或规定的签名输入格式。字段名称相同不表示语义相同，例如 security version 可能表示引导级防回滚值、应用版本或策略 generation；验证双方必须约定其作用域、单调规则和缺失处理。

\noindent\begin{minipage}{\textwidth}
\begin{lstlisting}[language=C,caption={远程证明响应的逻辑字段}]
struct attestation_evidence {
    uint32_t profile_version;
    uint8_t  device_identity[32];
    uint8_t  hardware_model[16];
    uint8_t  boot_measurement[32];
    uint8_t  runtime_measurement[32];
    uint8_t  challenge[32];
    uint64_t security_version;
    uint32_t lifecycle_state;
    uint32_t debug_state;
    uint64_t boot_count;
    uint8_t  signing_key_id[16];
    uint8_t  signature[64];
};
\end{lstlisting}
\end{minipage}

实际格式可使用 CBOR/COSE 等编码，逻辑结构不应直接映射为本地 C 结构体后跨网络发送。验证器必须检查类型、长度、重复字段、未知 critical claim、签名算法、证书链和 canonical encoding，拒绝“解析器看到一个值、签名器覆盖另一个值”的歧义。

Evidence 只陈述 Attester 能够观测并签名的事实。位置、租户、用户、业务健康或“设备无漏洞”等结论通常不应伪装成硬件 claim；它们应由验证方结合外部记录生成 Attestation Result。

\section{RATS 参与方与决策边界}

IETF RATS 架构区分 Attester、Verifier、Relying Party 和 Endorser\cite{rfc9334-rats}。Attester 产生 Evidence，Verifier 使用 endorsement、reference value 和策略进行评估，Relying Party 再根据 Attestation Result 决定是否允许接入、发放凭据或执行命令。

\begin{table}[H]
  \centering
  \caption{远程证明参与方及其输入输出}
  \begin{tabular}{p{27mm}p{38mm}p{42mm}p{29mm}}
    \hline
    参与方 & 可信输入 & 主要输出 & 不应承担的隐含职责 \\
    \hline
    Attester & 设备根、受保护状态、challenge & Evidence & 决定云端业务权限 \\
    Verifier & Evidence、endorsement、reference value、策略 & Attestation Result & 直接假定业务所有权 \\
    Relying Party & Result、账户与资源策略 & 接入或授权决定 & 重新解释原始签名字节 \\
    Endorser & 制造和供应链事实 & 已签名 endorsement & 决定现场设备当前健康 \\
    \hline
  \end{tabular}
\end{table}

Evidence 与 Attestation Result 必须在接口中使用不同类型。前者可能包含稳定硬件标识和详细 measurement，只应发送给受信验证器；后者可以是短期、最小披露的“允许/拒绝及原因类别”，供业务服务消费。把原始 Evidence 广播给所有微服务，会扩大隐私和解析攻击面。

证明成功只代表观测状态在给定时间满足指定策略，不等于设备所有应用无漏洞，也不等于调用者拥有目标资源。授权仍需绑定账户、租户、设备所有权、操作类型和会话状态。

\section{新鲜性与通道绑定}

签名正确的旧 Evidence 仍可能被重放。在线证明通常由 Verifier 或 Relying Party 提供不可预测 challenge，Attester 把它放入签名范围；验证器应一次性消费 challenge，并限制生成时间、接收时间和结果缓存时间。仅使用设备实时时钟不可靠，因为许多设备启动时尚未获得可信时间。

在离线或单向链路中，可以使用受保护启动计数、单调 evidence generation、短期时间窗口或预分配 challenge 集合，但每种方法都要说明状态丢失和同步恢复。普通文件中的递增整数不能抵抗整盘回放。

证明还应绑定正在建立的安全通道或待签发公钥。常见方式是把 TLS exporter、会话公钥摘要、CSR 公钥或协议 transcript 摘要放入 challenge/claim。否则中间人可以取得一台健康设备的 Evidence，再为另一条连接申请凭据。

缓存 Attestation Result 可以降低验证压力，但 cache key 至少应包含设备或伪名、策略版本、reference value generation、证据安全版本和通道绑定上下文。密钥吊销、基线撤回或紧急漏洞响应发生时，系统必须能使旧结果提前失效。

\section{Endorsement 与 Reference Value}

Endorsement 是由制造商、芯片厂或其他授权方签名的事实，例如某证明公钥属于某硬件型号、某启动级由特定供应链生成。Reference Value 是验证策略允许或已知的状态，例如某产品修订可接受的 Bootloader、TEE 和内核 measurement。两者都不是设备现场产生的 Evidence。

reference value 清单应包含产品型号、硬件修订、组件作用域、摘要算法、measurement、security version、有效期、状态和签名。分批发布期间通常需要同时接受旧稳定版、新候选版和恢复镜像，但每个值都应有启用条件、流量范围和撤回时间，不能永久积累“所有曾发布摘要”。

验证器应区分：

\begin{itemize}
  \item cryptographic invalid：签名、链或编码不成立；
  \item unknown：证据有效，但缺少型号 endorsement 或 reference value；
  \item non-compliant：事实可验证，但 security version、debug state 或 measurement 不满足策略；
  \item indeterminate：验证依赖暂时不可用或策略数据过期。
\end{itemize}

这些状态对应不同恢复路径。unknown 可能需要同步新基线，non-compliant 可能需要更新或隔离，invalid 通常应拒绝并告警；不能把所有错误都映射成“证书失败”或静默允许。

\section{证明密钥、隐私与规模}

证明密钥应与业务 TLS、软件更新和数据签名密钥分离。永久硬件根可派生 Initial Attestation Key，再按产品、租户或服务建立可轮换的 Attestation Key。DICE 架构展示了把设备秘密与逐层 measurement 组合、派生后续层身份的思路\cite{tcg-dice-architectures}；实际实现仍需固定派生上下文、证书格式和销毁边界。

向所有服务暴露永久唯一设备标识会形成跨域跟踪能力。根据业务需求，可以由受信验证器接收稳定身份，再向不同租户签发不可关联的伪名或短期结果；维修、转售和所有权转移时也应重新评估哪些标识继续可见。隐私策略必须说明谁能把伪名还原为硬件身份，以及审计和删除期限。

机群规模下，验证服务需要缓存厂商证书链、endorsement 和 reference value，而不是复用 challenge 或跨会话接受旧 Evidence。应对单设备、单租户和全局设置速率限制与并发上限，并监控签名失败、未知 measurement、结果延迟、吊销命中和时钟异常。

根或证明密钥疑似泄漏时，系统必须能撤销对应证书或 key ID，缩短 Result 有效期并隔离受影响批次。只按设备序列号拉黑不充分，因为攻击者可能复制序列号而继续使用被盗私钥。

\section{工作负载与服务身份}

硬件根密钥不应直接签署 MQTT 消息、业务订单或用户数据。更合理的模式是用设备证明引导短期工作负载凭据：验证器确认平台和目标工作负载状态，凭据服务再为本次生成且不可导出的公钥签发有限作用域证书或 token。

典型流程如下：

\begin{enumerate}
  \item 工作负载在受保护边界中生成会话密钥对；
  \item Relying Party 下发包含该公钥摘要的 challenge；
  \item Attester 返回绑定平台 measurement、工作负载 measurement 和 challenge 的 Evidence；
  \item Verifier 产生短期 Attestation Result；
  \item 凭据服务根据 Result、设备所有权和业务策略签发短期凭据；
  \item 服务端同时验证凭据、通道持钥证明和细粒度授权。
\end{enumerate}

这样可以在不暴露根密钥的情况下轮换业务凭据，并把授权范围限制到设备、服务、租户和有效期。证明不能替代进程隔离：若普通世界中的任意进程都能调用同一签名句柄，恶意进程仍可冒充健康工作负载。

\section{离线验证与失败策略}

完全离线的验证端需要预置并版本化信任锚、endorsement、reference value、策略和必要的吊销快照。验证包应带生成时间、有效期、适用产品、上一 generation 和签名，使现场工具能够判断“设备不合规”与“验证资料过旧”。

无可信实时时钟时，不能简单忽略证书有效期。可以使用安全版本、受保护计数、安装时记录的时间下界或带人工审批的维护窗口，但应把时钟不可信作为显式结果。恢复联网后要重新同步并复核离线期间接受的高风险操作。

失败策略取决于操作风险。读取非敏感遥测可在短期 indeterminate 状态下限权继续，发放长期密钥、开放调试或执行安全关键控制则应 fail closed。拒绝服务本身可能造成物理风险时，设备应进入预先定义的本地安全模式，而不是由云端证明故障直接切断全部控制。

诊断信息要足以定位证书链、nonce、measurement、策略版本或时间问题，但不应把私钥材料、完整稳定标识和可用于旁路的内部错误细节写入普通日志。

\section{生产注入与生命周期状态}

设备生命周期可划分为芯片制造、板级生产、产品装配、部署、维修和报废。每个阶段允许的调试、启动密钥和配置权限不同。

生产注入应使用受控工位、短期授权和可审计记录。私钥最好在设备内部生成，只导出公钥或证书请求；确需外部注入时，应使用安全通道并避免主机磁盘和日志保留明文。

生命周期状态宜由不可回退位、受保护计数或签名状态记录表示，例如 DEV、MFG、PROD、RMA 和 DECOMMISSIONED。状态必须进入 Evidence，并同时约束调试认证、启动密钥、测试命令和凭据签发。仅在普通文件中写入“production=true”没有安全意义。

每台设备的注入应是可恢复事务，关联芯片 UID、板卡序列号、硬件修订、公钥或 CSR、证书序列号、RPMB key 状态、镜像摘要、工位和最终锁定结果。中途失败的设备进入隔离队列，不能直接重跑并生成第二套身份。

从开发状态切换到量产锁定通常不可逆，必须在执行前验证安全启动、恢复镜像、设备证书和维修路径。锁定调试口却没有可用恢复路径，会把普通软件故障变成整机报废。

量产抽检不能只确认“证书存在”，还应执行一次完整 challenge-response、验证 measurement 与 lifecycle claim、模拟错误产品证书，并确认普通世界无法导出私钥。生产流程和组织审计的进一步要求见 \mbox{第~\ref{chap:security-lifecycle}~章}。

\section{密钥轮换、吊销、维修与报废}

根信任必须预留轮换机制。更新验证可以保留多个公钥槽位、签名授权链和安全版本计数；设备证书应支持续期和吊销。轮换流程必须在旧密钥仍有效时完成，避免设备因时间或网络故障永久失联。

不同密钥需要不同策略：不可变硬件秘密通常不能替换，只能停止信任对应设备；证明密钥可以通过上级身份重新背书；业务证书可以高频短期轮换；存储密钥轮换则需要数据重封装或版本化读取。把这些动作统一成一个“rotate key”接口会隐藏不可逆的数据后果。

轮换应采用有界双信任窗口：先发布新信任锚或授权链，确认目标机群已经接收，再启用新签名者，最后撤销旧密钥。验收需要覆盖长期离线设备、时钟错误、重复命令、轮换中掉电和服务端回滚。紧急泄漏处置可以缩短窗口，但不能跳过设备恢复路径。

维修更换主板时，原硬件身份不能简单复制到新板。系统应撤销旧设备、为新板建立身份，并通过所有权证明迁移租户与业务状态。若保留同一产品序列号，也必须在后台保存硬件身份 generation，避免新旧板同时有效。

设备报废需要撤销云端身份、清除用户数据和可删除密钥，并记录不可删除硬件身份的状态。恢复出厂设置和安全报废是不同操作。

\section{故障模式与恢复边界}

可信组件的错误路径应在产品状态机中有稳定含义，不能由上层把所有异常解释为“重试即可”。

\begin{table}[H]
  \centering
  \caption{典型可信机制故障与有界恢复}
  \begin{tabular}{p{28mm}p{38mm}p{42mm}p{27mm}}
    \hline
    故障 & 不安全处理 & 建议恢复边界 & 必要证据 \\
    \hline
    可信组件不可用 & 导出软件备用私钥 & 限权运行、维护或恢复系统 & 错误类别、原始码 \\
    RPMB 计数异常 & 忽略 generation & 冻结敏感写入并保留可读状态 & counter 与事务阶段 \\
    未知度量值 & 自动加入允许列表 & 隔离并同步签名基线 & 型号、组件与摘要 \\
    时间不可信 & 跳过全部有效期检查 & 使用受限离线策略并标记结果 & 时间源与下界 \\
    证明密钥吊销 & 只更换序列号 & 拒绝密钥并重新背书 & 证书链与吊销版本 \\
    \hline
  \end{tabular}
\end{table}

恢复系统本身也是启动和证明边界，必须有独立 measurement、最小命令集和认证路径。若正常系统无法证明就自动进入一个更宽松、可写任意分区的 recovery，会把拒绝服务转化为提权入口。

设备不应因验证服务离线而形成高速重试、RPMB 高频写入或 reboot loop。退避、最大尝试次数、离线宽限和人工维护入口都应有明确上限，并在恢复后上传最小审计摘要。

\section{验证、负向测试与故障注入}

可信功能测试不能止于一次正常签名。首先要从实现生成接口与状态清单，验证命令 ABI、claim schema、允许算法、生命周期转换、密钥属性和调试策略与设计一致。对 TEE、TPM 与证明解析器，应把边界检查纳入单元测试和模糊测试。

负向用例至少包括：

\begin{itemize}
  \item 截断、超长、重复、乱序、未知 critical claim 和非 canonical encoding；
  \item 错误签名算法、错误证书用途、过期链、被撤销 key ID 和错误硬件型号；
  \item 重放旧 nonce、并发复用 challenge、跨设备或跨 TLS 会话转移 Evidence；
  \item 合法签名但未知 measurement、较低 security version、开放调试和错误 lifecycle；
  \item 共享内存在检查后修改、整数溢出、取消与 TA 崩溃同时发生；
  \item 证明服务超时、reference value 更新中断和缓存提前失效。
\end{itemize}

故障注入要覆盖启动每一级复位、可信存储提交前后掉电、安全元件 BUSY/RESET、RPMB counter 不一致、world switch 中断、secondary CPU 上电和 suspend/resume。涉及 DMA 的平台应由非安全主设备主动扫描受保护地址，证明总线防火墙而不只是 CPU 页表生效。

验证器需要 differential test：不同语言或版本对同一 token 必须得到相同签名输入、claims 和决策；畸形 corpus 不得导致崩溃、无界内存或解析差异。密码操作还应测量并发下 P99/P99.9 延迟、失败率和节流行为，防止安全服务成为可被远程放大的容量瓶颈。

\section{数字化验收与评审清单}

项目应在测试计划中实例化 $T_{evidence}$、$T_{verify}$、$T_{fail}$、$N_{power}$、$T_{revoke}$ 和计数器寿命预算 $B_{counter}$。这些符号不能留在最终发布门中；每个硬件修订和安全 profile 都要给出数值、负载、温度、网络和置信要求。

\begin{table}[H]
  \centering
  \caption{可信执行与证明验收矩阵}
  \begin{tabular}{p{29mm}p{39mm}p{39mm}p{27mm}}
    \hline
    验收项 & 量化门槛 & 测试条件 & 留存证据 \\
    \hline
    Evidence 生成 & P99 $\leq T_{evidence}$，超时有界 & 冷/热机、并发、低电压与高温 & 分位数与错误分类 \\
    在线验证 & P99 $\leq T_{verify}$，重放 nonce 全拒绝 & 正常、吊销、未知基线与缓存失效 & token、策略版本、判定 \\
    错误收敛 & 最迟 $T_{fail}$ 进入规定状态 & TEE、网络、时钟和安全元件故障 & 状态轨迹与审计码 \\
    掉电一致性 & $N_{power}$ 次切电均不接受旧状态 & 提交各阶段随机切电 & seed、波形与恢复根 \\
    吊销传播 & 全部验证器在 $T_{revoke}$ 内拒绝 & 缓存、分区和回滚场景 & 配置版本与命中率 \\
    计数器寿命 & 累计写入 $< B_{counter}$ 且有裕量 & 正常、异常重试和维修流程 & 写入模型与实测计数 \\
    \hline
  \end{tabular}
\end{table}

发布评审还应确认：

\begin{itemize}
  \item 根身份是否来自不可由普通软件修改的信任源？
  \item 各用途密钥是否独立派生，context 是否版本化并具有迁移路径？
  \item TCB 是否有清单，TEE 接口是否把所有共享内存视为不可信输入？
  \item CPU、DMA、调试器、中断、缓存和复位是否处于一致隔离边界？
  \item measured boot 是否覆盖 DTB、命令行、initramfs、恢复镜像和协处理器固件？
  \item Evidence、Endorsement、Reference Value 与 Attestation Result 是否使用不同类型和权限？
  \item challenge 是否一次性消费，证明是否绑定会话或待签发公钥？
  \item reference value 是否签名、版本化，并具有灰度启用与紧急撤回路径？
  \item 验证器不可用、时钟不可信和计数器耗尽是否进入有界安全状态？
  \item 量产锁定、离线验证、维修换板、密钥轮换、吊销和报废是否完成演练？
  \item 全部验收门槛是否已填入具体数值，并关联原始日志、波形和构建版本？
\end{itemize}
